\documentclass[fontsize=16pt]{mybibreport}


\graphicspath{{../../../assets/images/}}
\setheaderlogo{mnr-green}

\title{Predigt}
\subtitle{Vom Ölberg nach Jerusalem}
\author{Schmid Lothar}
\date{5/7/2026}

\begin{document}

    \mytitlepage

    \tableofcontents

    \chapter{Gott}

    \section{Einführung}
    6 Tage vor den Passahfest war Jesus bei Lazarus und seinen Schwestern zu Besuch. Jesus hat Lazarus von den Totenauferweckt und diese Familie gehörte zu seinen Freunden. Lazarus und seine Schwestern wohnten in Betianien, welches ca. 3km im Osten von Jerusalem nahe dem Ölberg liegt.
    Es gibt aber noch ein zweites Betanien, das in \begin{bibelbox}{SCHL}{Joh}{1:28}
    
    \end{bibelbox}
    Oft wird dieses Betanien auch als Bethabum übersetzt um eine Verwechslung mit unserem Betaninen, in dem Jesus bei lazarus zu Besuch war. Bei diesem Besuch hat Maria Jesus die Füsse mit kostbaren Öl gesalbt. Worauf Judas der Verräter, die Frau zurechtwies. Jesus antwortete darauf:
    \begin{bibelbox}{SCHL}{Joh}{12:7}
    
    \end{bibelbox}
    Jesus hat bei Lazarus übernachtet im Vers 12 lesen wir:
    \begin{bibelbox}{SCHL}{Joh}{12:12}
    
    \end{bibelbox}
    Jesus ist also am Morgen Richtung Jerusalem aufgebrochen. Das Haus von Lazarus, wurde aber regelrecht von Menschen belagert, die von dem Wunder der Totenauferweckung gehört haben. Sie wollten diesen lebenden Toten anschauen und passte gerade gut das der Wunderheiler auch gerade vor Ort war.

    Als Jesus uns seine Jünger sich auf den Weg nach Jerusalem machten, begleitete ihn die Menge. Unterwegs wies Jesus zwei Jünger an, nach Jerusalem zu gehen und dort ein Füllen zu holen, auf dem er beabsichtigte in die Stadt zu reiten. Das lesen wir in der Parallelstellen von Matthäus, Markus und Lukas.
    \begin{bibelbox}{SCHL}{Matt}{21:1-3}
    
    \end{bibelbox}
    Die Menschen um ihn herum legten Palmzweige aus, die Jünger legten ihre Kleider auf den Esel und auf den Weg und alle preisen den Herrn:
    \begin{bibelbox}{SCHL}{Joh}{12:13-15}
    
    \end{bibelbox}
    Ich habe im Internet einen Gottesdienst von Amerika gesehen. Es war eine riesige Halle voll von Menschen. Der Pastor auf der Bühne schrie nur immer: \enquote{Amen, Amen...} nur dieses eine Wort, und die Menge schrie begeistert mit. Es ist einfach Massen mitzureissen und zu begeistern, sieht man jetzt auch an der Fussball-WM wie die Leute in den Stadien im Chor singen.

    Ich denke viele in diesem Gottesdienst waren ergriffen und lobten wirklich Gott. Wir können mit dem Finger auf diese Menschen zeigen, aber wenn man mit einem Finger auf den Gegenüber zeigt, zeigen immer vier Finger gegen sich selber.

    So zum Beispiel das Lied \enquote{Grosser Gott wir loben dich}, ein wunderschönes Lied, wird oft mit unglaublicher inbrunst gesungen. Ich kann mich gut an die Ostermessen hier in Naters erinnern, wo dieses Lied in einer vollen Kirche mit Orge und Chor gesungen wird, einfach gewaltig. Und ja, ich bin überzeugt, bei vielen kommt es in dem Moment von Herzen. Wie bei der Menschenmenge die Jesus begleiteten.

    In der Bibel lesen wir von mehreren solchen Momenten im Alten Testatment, bei der Einweihung des Tempels zum Beispiel \bibleverse{IKoen}(10:10). Oder als Nehemia nach dem Mauerbau und der Schriftauslegung durch Esra und die Priester das Laubhüttenfest wieder aktivierte. Oder beim König .... der die Schrift im Tempel gefunden hat und darauf das Passahfest so feierte wie seit den Richtern nicht mehr.

    Gott hat uns Begeisterungsfähig geschaffen und wir sollen das auch sein, vor allem können wir begeistert von seiner Gnade und wohltat sein. Aber was passiert, wenn wir solche Gottesdienste, uns bis ins innerster begeistern verlassen? Wie lange hält sich diese Hochstimmung an? Ich weiss noch früher in Birgisch, nach dem Amt, ging es rüber ins Restaurant zum Apero. Nach ein paar Ballon Weisswein oder Bier hat sich die Begeistrung schnell ins Weltliche gedreht. Die Frauen gingen früher nach Hause an den Herd und die Männer später, mancher mehr manche weniger betrunken. Es muss ja nicht Restaurant und trinken sein. Nach den Gottesdiensten geht viel tratsch und klatsch von einer Person zu anderen über.

    Am Montag wenn dann die hektische Woche wieder anfängt, ist dann von dieser sonntäglichen Begeisterung nicht mehr viel zu spüren. So leben wir oft einfach von Sonntag zu Sonntag. Es ist aber wichtig, dass wir diese Begeisterung wärend der Woche nicht versiegen lassen, sondern uns durch Gebet und das Lesen im Wort immer wieder neu von Jesus begeistern lassen.

    Es war also ein grosser Zug der sich vom Ölberg Richtung Jerusalem aufmachte. Lobgesängen und Palmzweige und Jubel. Mitten drin Jesus auf dem Füllen. Jesus wusste genau, dass viele von denen die ihn hier jetzt zu Jubeln und Lobpreisen, in ein paar Tagen, \enquote{Kreuzige ihn!} schreien werden.

    Das ist es, was wir nie vergessen sollten. Gott weiss alles, er weiss alles schon im Voraus. Er weiss wie wir uns nach dem Lobsang am Sonntag wärend der Woche verhalten werden.

    Kennt ihr die netten freundlichen Gäste, die so über euer Essen schwärmen und wie toll doch alles war. Später von anderen Freunden wird euch zugetragen, dass die Gäste über euch wegen dem Essen gelästert haben. Was denkt ihr über solche Gäste? Würdet ihr sie wieder einladen?

    Gott ist da ganz anders. Als Adam und Eva Gott hintergangen haben, hat Gott nach Adam gesucht.
    \begin{bibelbox}{SCHL}{IMos}{3:9}
    
    \end{bibelbox}
    Auch das Volk Israel hat er immer wieder gesucht obwohl sie von ihm angefallen sind. 
    \begin{bibelbox}{SCHL}{Jes}{1:1}
        
    \end{bibelbox}
    So auch hier, Jesus wusste, dass viele in diesem Volk die ihn hier jetzt loben und preisen später ihn schreien werden: \enquote{Kreuzige ihn}. Auf wusste er, dass die Jünger ihn in der schwersten Stunde verlassen werden, dass er von Petrus verleugnet wird. Aber Jesus ist seinen den von Gott ersonnen Weg weiter gegangen. Er ist trotz der Ablehnung der Menschen ans Kreuz für uns gegangen. Das aus reiner Liebe:
    \begin{bibelbox}{SCHL}{Joh}{3:16}
    
    \end{bibelbox}

    In Vers 19 das die Pharisäer langsam in Panik kommen, als sie diese Menge sahen, die so begeistern den Heiler und Wanderprediger feiern. Sie bemerken, dass diese Bewegung ihnen langsam ihrer Kontrolle entgleitet und aus dem Ruder läuft. Als sie dann Jesus ans Kreuz geliefert hatten, dachten sie jetzt hätten sie es Griff, jetzt wird die Bewegung in sich zusammenfallen, wie viele andere. Aber dem war nicht so. An Pfingsten ging es erst richtig los und zwar so, dass Gamaliel ein führender Pharisäer sagen musste
    \begin{bibelbox}{SCHL}{Apg}{5:38}
    
    \end{bibelbox}

    Heute wissen wir, dass es Gottes Wille ist. Er hat Paulus berufen und die Gemeinde wurde aufgebaut. Israel wurde für einen Moment auf die Seite gestellt und die Gemeinde der Leib Christi wurde aufgebaut.
    \begin{bibelbox}{SCHL}{Rom}{11:1}
    
    \end{bibelbox}
    Die ganze Welt wird von der christlichen Bewgung geprägt. Die ganze Welt ist durchdrungen vom Christentum, überall wird versucht dieses zu unterdrücken, durch Gewalt oder wie hier bei uns durch Bibelkritik und Wissenschaft. Es heisst, das in Nordkorea mehr Christen im KZ sitzen, wie es hier in der Schweiz wiedergeborene Christen gibt.
    Wir wissen aber, dass es bestehen bleibt, bis wir entrückt werden.

    Nicht nur, dass Jesus verkündigt, dass er auf dem Ölberg wiederkommt, 
    \begin{bibelbox}{SCHL}{Apg}{1:11}
    
    \end{bibelbox}
    sondern, dass er sogar in Macht und Herrlickeit wiederkommt und das Regiment hier auf der Erde übernimmt. 
    \begin{bibelbox}{SCHL}{Offb}{1:1}
    
    \end{bibelbox}
    Da wird wieder ein jubel vor den Toren Jerusalem sein. Jerusalem dann zwar in Trümmer liegen, aber eine 1000jährige Zukunft in Frieden liegt vor den Menschen.

    Zurück zu unserer Geschichte. Jesus ist in dieser Woche den Weg mehrmals gegangen.
    \begin{bibelbox}{SCHL}{Lk}{21:37-38}
    
    \end{bibelbox}
    Es ist gut möglich, dass Jesus mit seinen Jünger bei Lazarus übernachtet haben. 
    \begin{bibelbox}{SCHL}{Mk}{11:11}
    
    \end{bibelbox}
    Auf einem von diesen Rückwegen von Jerusalem nach dem Ölberg hielt Jesus den Jüngern seine berühmte Ölberbergrede, die in \bibleverse{Matt}(24:)ff 

    Es muss gewaltig gewesen sein, wenn die kleine Gruppe auf dem Ölberg im Kreis sitzen mit dem Blick auf das gewaltige Bauwerk den Tempel. Heute sehen wir den ..-dom mit der goldenen Kuppel. Damals war es ein Bauwerk, dass Joseph Flavius als das grösste und schönste Bauwerk zu seiner Zeit beschreibt.

    Wie Jesus prophezeit hat, findet man heute nichts mehr von diesem Tempel. Er wurde 70 und 130 \nChr total zerstört. Die Klagemauer an der heute die Juden beten ist ein Teil der Stadtmauer und kein Teil von dem Tempel.

    So wie diese Vorhersage von Jesus in Erfüllung ging, gehen auch alle anderen seine Vorhersagen in Erfüllung. So, \zB dass er für alle sichtbar auf dem Ölberg wiederkommt.
    \begin{bibelbox}{SCHL}{Offb}{1:1}
    
    \end{bibelbox}
    Wann die Entrückung stattfindet wissen wir nicht. Es so wie es in den Lied, das wir oft singen heisst, \enquote{Wir wissen nicht den Tag noch die Stunde}. Aber Jesus kommt wieder und wir mit ihm als seine Braut
    \begin{bibelbox}{SCHL}{Dan}{7:13-14}
    
    \end{bibelbox}
    \begin{bibelbox}{SCHL}{Offb}{19:11.14}
    
    \end{bibelbox}
    So lange, das nicht eintrifft sind wir in der Gnadenzeit.

    Am letzten dieser Woche, hat Jesus mit den Jüngern das Abenmahl gefeiert. An diesem Abend hat Jesus den Jüngern gesagt, tut dies zu meinem Gedächtnis. Später hat Jesus Paulus aufgetragen, dieses Mal so aufzuschreiben, dass es nicht vergessen wird. Darum wollen wir auch heute gemeinsam das Mal zum an das Opfer welches er uns ein für alle mal für uns am Kreuz gemacht hat.
    \begin{bibelbox}{SCHL}{IKor}{11:26}
    
    \end{bibelbox}

\end{document}